% ----------------------------------------------------------
% Introdução (exemplo de capítulo sem numeração, mas presente no Sumário)
% ----------------------------------------------------------
%\chapter*[Introdução]{Introdução}
\chapter{Contextualização}
%\addcontentsline{toc}{chapter}{INTRODUÇÃO}
% ----------------------------------------------------------

\section{Introdução}
O ensino de ciências, e em particular o de Física, é historicamente pautado no equilíbrio entre a formalização matemática e a experimentação prática \cite{optics}. A atividade experimental não é apenas um complemento à teoria; ela funciona como um ambiente didático que permite ao estudante manipular variáveis, confrontar modelos ideais com ruídos do mundo real e reestruturar seus esquemas mentais \cite{raviola}. No entanto, a implementação eficaz de laboratórios didáticos enfrenta barreiras estruturais severas, especialmente em instituições públicas, onde o alto custo de aquisição e manutenção de sistemas proprietários de Aquisição de Dados (DAQ) limita a modernização e a escalabilidade das práticas. 

A evolução das Tecnologias Digitais de Informação e Comunicação (TDIC) trouxe novas possibilidades, como os Laboratórios de Acesso Remoto (LAR) e as experiências híbridas, que utilizam aparatos reais controlados via rede e internet \cite{lms}. Esse movimento ganhou urgência sem precedentes com a pandemia de COVID-19, que forçou instituições ao redor do mundo a reinventar seus métodos de ensino para garantir a segurança de alunos e docentes sem comprometer os resultados de aprendizagem \cite{create, transforming}. A arquitetura proposta neste trabalho insere-se nessa fronteira tecnológica, buscando não apenas medir fenômenos físicos, mas prover uma infraestrutura orquestrada que suporte a transição do laboratório isolado para um ambiente conectado e centralizado.

\subsection{Contextualização do Tema}
A literatura atual destaca que a eficácia pedagógica de um laboratório está intrinsecamente ligada à capacidade do aluno de agir como um cientista: formulando hipóteses e coletando dados primários em ambientes reais, em vez de apenas observar simulações virtuais que simplificam a complexidade da natureza. Para viabilizar essa prática em larga escala, surgiram abordagens baseadas em hardware de código aberto, como Arduino e ESP32, que democratizaram o acesso à instrumentação científica ao reduzir custos de implementação em relação aos kits comerciais da Pasco ou Vernier.

No cenário da instrumentação científica educacional, o mercado global é amplamente dominado por ecossistemas proprietários de empresas como Pasco Scientific, Vernier (LabQuest) e TeachSpin \cite{kraftmakher}, com forte representação de equipamentos DAQ (\textit{Data Acquisition Systems}). Embora essas soluções ofereçam alta precisão e recursos avançados de processamento, sua adoção em larga escala esbarra em barreiras financeiras e arquiteturais intransponíveis para a maioria das instituições públicas de países em desenvolvimento \cite{vidal}. Estima-se que sistemas comerciais completos para experimentos específicos possam ultrapassar o custo de milhares de dólares por bancada \cite{mishonov}, um valor proibitivo que acentua a desigualdade no acesso à experimentação de qualidade. Além do alto investimento, tais sistemas são frequentemente projetados como unidades isoladas que exigem um computador dedicado por experimento para a visualização dos dados via conexão física (USB), resultando em uma infraestrutura rígida e com alto índice de ociosidade de hardware durante os períodos de discussão teórica. Esta configuração impede que o laboratório seja gerenciado como uma rede integrada, dificultando a supervisão docente e a transição para modelos de ensino mais modernos e dinâmicos.

Entretanto, a simples substituição de DAQs caros por microcontroladores não resolve o problema de gerenciamento do ambiente laboratorial. Frequentemente, essas soluções operam de forma fragmentada: cada bancada exige um computador dedicado para a visualização dos dados via cabo USB, ou depende de smartphones pessoais dos alunos. Embora o uso de smartphones (BYOD - \textit{Bring Your Own Device}) ofereça portabilidade, ele introduz desafios: a falta de uma rede de dados integrada impede que o docente monitore simultaneamente o progresso de múltiplos grupos, sobrecarregando a gestão da aula e dificultando intervenções em tempo real.

A tendência mais recente na instrumentação educacional foca no desenvolvimento de frameworks escaláveis, como o FREE (\textit{Framework for Remote Experiments in Education}), que defendem a separação entre a interface de usuário (UI) e o controle do aparato físico via protocolos web modernos. Essa abordagem permite que o laboratório funcione de forma autônoma, onde múltiplos nós sensores (IoT) transmitem telemetria sem fio para um ponto central de persistência e visualização.

\subsection{Apresentação e Delimitação do Problema}
O problema central identificado é o isolamento arquitetural das bancadas experimentais e a consequente ociosidade de hardware. Em uma configuração tradicional, o hardware de coleta permanece inativo durante os longos períodos destinados à discussão teórica e montagem, e sua dependência de conexões físicas individuais impede a integração dos dados. Este cenário gera uma lacuna tecnológica: como equilibrar a economia de recursos de sistemas de baixo custo com a robustez técnica necessária para uma orquestração centralizada?

A delimitação deste trabalho foca no desenvolvimento de uma infraestrutura IoT escalável, baseada em microcontroladores ESP32 e um servidor Raspberry Pi, que utiliza o protocolo Wi-Fi para eliminar o emaranhado de cabos e centralizar a coleta de dados de múltiplos experimentos. O escopo abrange a demonstração técnica da arquitetura através de uma Prova de Conceito (PoC) com o experimento de Carga e Descarga de Capacitor, por meio do desenvolvimento tecnológico, focando em proporcionar a viabilidade da orquestração dos dados e trazer o diferencial da telemetria sem fio.

\section{Justificativa e Contribuições}
A justificativa para este estudo reside na necessidade de fornecer soluções de Sistemas de Informação que apoiem a modernização dos laboratórios de ensino de forma economicamente viável. Diferente de trabalhos que focam apenas no uso de um único sensor, esta pesquisa justifica-se pela criação de uma camada de orquestração, permitindo que o docente gerencie o laboratório como uma rede de dispositivos inteligentes, ccom seus sensores e montagens associadas cada. A relevância tecnológica manifesta-se na proposição de uma arquitetura que reduz drasticamente a necessidade de hardware duplicado (como computadores individuais por bancada), otimizando o investimento institucional e permitindo que o foco pedagógico retorne para a análise dos dados e para a interação professor-aluno.

Do ponto de vista científico, o trabalho contribui para a literatura de experimentação remota e híbrida ao demonstrar como protocolos leves de comunicação sem fio podem ser aplicados para garantir a sincronização de dados experimentais reais \cite{lms}. Ao integrar ferramentas de IoT a interfaces web amigáveis, o projeto serve de alicerce para futuros laboratórios remotos que operam 24 horas por dia, por exemplo, permitindo que o conhecimento físico seja acessado de forma democrática e independente do espaço físico da universidade \cite{termometria}.

\subsection{Pergunta de Pesquisa}
Diante das limitações de custo, escalabilidade e monitoramento nos laboratórios de Física, este trabalho busca responder: Como projetar e implementar uma infraestrutura de baixo custo, baseada em comunicação sem fio e arquitetura IoT, que permita a centralização, persistência e visualização de telemetria experimental que possibilite múltiplas bancadas em um ambiente educacional?

\section{Objetivos}
A definição de objetivos claros é fundamental para delimitar o escopo da pesquisa e estabelecer as metas técnicas que guiarão o desenvolvimento do artefato. Nesta seção, apresentam-se o objetivo geral e os objetivos específicos, estruturados sob a ótica da resolução do problema de isolamento e ociosidade de hardware em laboratórios de Física.

\subsection{Objetivo Geral}
Projetar, implementar e demonstrar a funcionalidade de uma arquitetura de monitoramento centralizado baseada em Internet das Coisas (IoT), utilizando microcontroladore ESP32 e um servidor Raspberry Pi via rede Wi-Fi, visando a orquestração e a visualização em tempo real de um experimento didático de Física em um ambiente laboratorial integrado.

\subsection{Objetivos Específicos}
Os objetivos específicos detalham as etapas técnicas e investigativas necessárias para a concretização da arquitetura proposta.

\begin{enumerate}
    \item Implementar uma Prova de Conceito (PoC) utilizando o experimento de Carga e Descarga de Capacitor (Circuito RC), por ser um arranjo de alta relevância didática e facilidade de digitalização, servindo como base para avaliar o funcionamento da telemetria digital, sendo programado o firmware em um microcontrolador ESP32.
    \item Implementar a camada web em código Python, culminando em um servidor conectado a um banco de dados, com interface gráfica para o usuário final.
    \item Demonstrar a viabilidade técnica da orquestração de dados, documentando leituras experimentais dispostas na interface web, consolidando o artefato como uma infraestrutura funcional para laboratórios híbridos e remotos.
\end{enumerate}
